% Part: first-order-logic
% Chapter: lindstrom
% Section: introduction

\documentclass[../../../include/open-logic-section]{subfiles}

\begin{document}

\olfileid[hi]{mod}{lin}{int}
\section{परिचय}

इस अध्याय में हमारा लक्ष्य लिंडस्ट्रॉम (Lindstr\"om) के उस अभिलक्षणन को
सिद्ध करना है जिसके अनुसार प्रथम-कोटि तर्कशास्त्र वह अधिकतम तर्क है जिसके
लिए (कुछ अतिरिक्त प्रतिबंध दिए होने
पर) संहतता और अवरोही लोवेनहाइम--स्कोलेम (L\"owenheim--Skolem) प्रमेय मान्य हैं
(\olref[fol][com][com]{thm:compactness} और
\olref[fol][com][dls]{thm:downward-ls})। पहले हमें उन तर्कों के सामान्य
वर्ग का अधिक व्यापक निरूपण चाहिए जिन पर यह प्रमेय लागू होता है। हम अपने को
\emph{संबंधपरक} भाषाओं तक सीमित रखेंगे, अर्थात् ऐसी भाषाओं तक जिनमें केवल
!!{predicate}s और व्यष्टि-अचर हों, पर कोई !!{function} न हो।

\end{document}
